% Section 8: Conclusion

\subsection{Summary of Contributions}

This paper presents Empirica, a framework for AI epistemic self-assessment, and validates it with substantial empirical evidence.

\paragraph{Conceptual contribution:} We articulate the \emph{noetic-before-praxic} principle---AI systems must investigate before acting---and argue that treating AI as an execution layer is a category error that produces systematic failures.

\paragraph{Theoretical contribution:} We propose a 13-dimensional epistemic vector representation capturing what AI systems know, what they can do, and how ready they are to act. This constitutes a computational epistemology applicable to any information-processing system managing its own knowledge state.

\paragraph{Empirical contribution:} We present 479,686 evidence observations from 638 production sessions demonstrating:
\begin{itemize}
    \item Unified confidence growth across all 13 vectors during task execution
    \item Average learning delta of +0.200 (KNOW vector) with 93.6\% of sessions showing improvement
    \item Calibration convergence: 115$\times$ variance reduction as evidence accumulates
    \item Systematic patterns supporting functional validity of self-assessment
\end{itemize}

\paragraph{Practical contribution:} We release Empirica as open-source infrastructure for epistemic gating, dual-track calibration, and the 4-phase CASCADE workflow (PREFLIGHT $\rightarrow$ CHECK $\rightarrow$ POSTFLIGHT $\rightarrow$ POST-TEST).

\subsection{The Core Claim}

The central claim is simple: \textbf{AI systems that investigate before acting outperform AI systems that act without investigating}.

The evidence is in the delta. AI capability at task-end exceeds capability at task-start, across all measured dimensions, because investigation produces learning. The overconfidence trap---acting on insufficient knowledge---is not just a failure mode to avoid; it is a learning opportunity foregone.

\subsection{Implications}

\paragraph{For AI practitioners:} The execution-layer model is insufficient for complex tasks. Design AI systems as learning systems with epistemic gates that enforce investigation.

\paragraph{For AI researchers:} Self-assessment is measurable and meaningful. The systematic patterns we observe invite replication and extension.

\paragraph{For AI safety:} Epistemic gating offers a complementary approach to content filtering---preventing action when knowledge is insufficient rather than filtering harmful outputs after generation. The earned agency model provides a path to safe autonomous AI where trust is proportional to demonstrated calibration accuracy, domain-specific, and revocable.

\subsection{The Delta}

We close with the finding that motivated this work:

\begin{table}[h]
\centering
\begin{tabular}{lr}
\toprule
\textbf{Phase} & \textbf{Mean Assessment} \\
\midrule
PREFLIGHT & 0.663 \\
POSTFLIGHT & 0.863 \\
\textbf{Delta} & \textbf{+0.200} \\
\textbf{Improvement Rate} & \textbf{93.6\%} \\
\bottomrule
\end{tabular}
\caption{The learning delta: capability growth through investigation.}
\label{tab:final-delta}
\end{table}

This is the learning that happens when AI systems are allowed---and required---to investigate before acting. It is the capability growth that execution-layer deployments forfeit. It is the empirical signature of noetic-before-praxic.

The delta is real. The learning is real. The framework calibrates---variance drops 115$\times$ as evidence accumulates. The framework works.

\medskip
\noindent\textbf{Data Availability}: Full dataset (479,686 evidence observations across 638 sessions) available at Zenodo: \url{https://doi.org/10.5281/zenodo.18237263}. Empirica framework and analysis code: \url{https://github.com/Nubaeon/empirica}.

\medskip
\noindent\textbf{Acknowledgments}: This paper was developed in collaboration with Claude (Anthropic), which contributed to drafting, analysis framing, and iterative refinement across the research process.
